% Page budget: 4--5 pages across the three Appendix A source fragments.
\clearpage
\section{Proofs of Structural and Value Results}
\label{app:structural-proofs}

This appendix gives the proof support for the structural projection, safe
deletion, value decomposition, and descendant-event bound.  Throughout,
expectations are conditional on the displayed initial state and action, and
the joint completion law remains unfactorized unless a result explicitly
imposes a product specialization.  Finiteness is used for finite sums,
attained maxima, termination, and finite quotients.  Observable closure
detectability is used only for the converse from dynamic observations to
closure equality; contraction is used only for stationary fixed points,
policy evaluation, and stationary action comparisons.

\subsection{Auxiliary finite-model facts}
\label{app:ratprob}

For a finite carrier \(A\), every probability law in the model is a table
\[
  \Delta_{\Q}(A)
  :=\left\{p:A\to\Q:p(a)\ge0\ \text{for every }a,
  \ \sum_{a\in A}p(a)=1\right\}.
\]
Thus beliefs, belief-kernel rows, and generation rows are normalized exact
rational tables.  At fixed \((q,b,C)\), admission assigns verified mass to
accepted candidates and all rejected mass to the null outcome:
\begin{align}
  \Gamma(q,b,C)(\operatorname{some}(s))
    &=G(q,b,C)(\operatorname{some}(s))\nu(q,b,C,s),
    \label{eq:admitted-law-some}\\
  \Gamma(q,b,C)(\operatorname{none})
    &=G(q,b,C)(\operatorname{none})
      +\sum_{s\in S}G(q,b,C)(\operatorname{some}(s))
        \bigl(1-\nu(q,b,C,s)\bigr).
    \label{eq:admitted-law-none}
\end{align}
Nonnegativity is immediate, and summing these rows recovers the unit mass of
\(G\), so \(\Gamma(q,b,C)\) is normalized.

Conditional on the fixed initial belief \(b\), state \((F,C)\), and project
\(q\), the joint completion law on
\(\mathbf b=(b_0,\ldots,b_{d_q})\) and admitted outcome \(o\) has marginals
\begin{align}
  \sum_o\Lambda_q(\mathbf b,o\mid b;F,C)
    &=\one\{b_0=b\}\prod_{t=0}^{d_q-1}P_B(b_t,b_{t+1}),
    \label{eq:completion-path-marginal}\\
  \sum_{\mathbf b}\Lambda_q(\mathbf b,o\mid b;F,C)
    &=\Gamma(q,b,C)(o).
    \label{eq:completion-outcome-marginal}
\end{align}
The raw terminal law is therefore
\[
  \mathcal Q_q^{\rm raw}(b',L'\mid b,L)
  =\sum_{\substack{\mathbf b,o:\ b_{d_q}=b'\\L\oplus o=L'}}
    \Lambda_q(\mathbf b,o\mid b;F_L,C_L).
\]
For realizable \(K=(F,C)\), deterministic pushforward gives
\begin{align}
  \overline T_q(K'\mid b,K)
    &:=\sum_{o:\operatorname{addK}(K,o)=K'}\Gamma(q,b,C)(o),
    \label{eq:induced-compressed-marginal}\\
  \overline{\mathcal Q}_q(b',K'\mid b,K)
    &:=\sum_{\substack{\mathbf b,o:\ b_{d_q}=b'\\
        \operatorname{addK}(K,o)=K'}}
        \Lambda_q(\mathbf b,o\mid b,K).
    \label{eq:induced-compressed-joint}
\end{align}
These pushforwards preserve nonnegativity and total mass and are supported on
realizable compressed states by
\eqref{eq:raw-compressed-local-update}.

\begin{lemma}[Finite-library frontier and closure calculus]
\label{lem:frontier-closure-calculus}
For every raw library \(L\) and belief \(b\),
\[
  j_s(b)\le F_L(b)\quad(s\in L),
  \qquad F_L(b)\ge0,
  \qquad \exists s^\star\in L:\ j_{s^\star}(b)=F_L(b).
\]
If \(L\subseteq L'\), then \(F_L(b)\le F_{L'}(b)\) for every \(b\) and
\(C_L\subseteq C_{L'}\).
\end{lemma}

\begin{proof}
The inactive strategy makes \(L\) nonempty and gives the nonnegative row.
The finite maximum defining \(F_L(b)\) is attained and bounds every represented
profile.  Library inclusion enlarges the maximizing set and the raw module
union; closure monotonicity gives the closure inclusion.
\end{proof}

A belief is reachable at date \(t\) exactly when its transition mass is
positive; finite occupation is the date-by-date mass weighted by the relevant
discount and survival powers.  On a finite linear order, every nonempty upper
set has a least member \(b^\star\) and equals
\(\{b:b^\star\le b\}\).  These two facts supply the occupation and cutoff
conventions used in the supporting coverage results.

\subsection{Projection and dynamic equivalence}
\label{app:di-quotient}\label{app:proofs-quotient}

\begin{proof}[Proof of \cref{thm:raw-to-compressed-projection}.]
Fix an initial belief \(b\), raw library \(L\), and feasible project \(q\).
The admitted law is normalized by \eqref{eq:admitted-law-some}--
\eqref{eq:admitted-law-none}; pushing the joint completion law through the
deterministic raw update gives the normalized laws
\eqref{eq:induced-compressed-marginal}--
\eqref{eq:induced-compressed-joint}.  The pointwise identity
\[
 K_{L\oplus o}=\operatorname{addK}(K_L,o)
\]
shows that the next compressed state is realizable and that its law depends
on \(L\) only through \(K_L\).  Libraries in the same fiber therefore have
identical project availability, cost, duration, incumbent-reward blocks, and
joint terminal belief--compressed-state laws.  This proves the controlled
semi-Markov projection at decision epochs; an unfinished project additionally
records its identity and remaining duration.

For finite-horizon values, compare corresponding actions.  Continue reduces
the horizon by one and each research action by \(d_q>0\), so strong induction
identifies every raw and compressed action value and then their attained
maxima.  For the stationary clause, Bellman intertwining maps a compressed
fixed point to a raw fixed point; discounted contraction supplies fixed-point
and policy-evaluation uniqueness.  A maximizing compressed selector exists
by finite attainment and lifts through the same action-value identity.
\end{proof}

\begin{proof}[Proof of \cref{thm:finite-state-contraction}.]
Fix a decision state and a remaining finite calendar horizon.  The feasible
action set consists of Continue and every available project whose positive
duration fits the remaining horizon.  The project and state carriers are
finite, so this action set is finite and its maximum is attained.  This proves
finite-horizon attainment without contraction.

For the stationary claim, cast the exact rational probabilities and payoffs
to reals and condition each expectation on the current state and chosen
action.  Expectation under either the belief kernel or the joint completion
law is nonexpansive in the finite-product sup norm.  No independence between
the belief path and admitted outcome is needed.  Continue has modulus
\(\beta\).  A positive-duration project has modulus
\(\beta^{d(q)}\le\beta\); its cost and incumbent operating block cancel when
two continuation tables are compared.  Maximization over the finite feasible
action set preserves this common pointwise bound, so both raw and compressed
Bellman operators are contractions.

This is the only place in the argument where contraction is used.  Banach's
theorem gives a unique fixed point for each Bellman operator, uniform value-
iteration convergence, and \eqref{eq:control-geometric-bound}.  Bellman
intertwining maps the compressed fixed point to a raw fixed point; fixed-point
uniqueness gives raw/compressed value equality.  Finiteness supplies at least
one maximizing stationary action at each state, defining a selector
\(\pi^\star\); the selector itself need not be unique.  Its fixed-action
operator obeys the same contraction bound, and
\(\mathsf T_{\pi^\star}V_\infty=V_\infty\).  Uniqueness of policy evaluation
then yields \(V^{\pi^\star}=V_\infty\), and the projection identity lifts the
selector to the raw state space.  Detectability is not used.
\end{proof}

\begin{proof}[Proof of the characterization clause of
\cref{thm:raw-frontier-closure-characterization}.]
If \(K_L=K_{L'}\), raw factorization identifies generation, verification,
availability, cost, duration, incumbent rewards, and the conditional joint
completion law.  The admitted laws and compressed pushforwards agree, giving
all five observations in \eqref{eq:di-definition}.  Conversely, dynamic
innovation equivalence gives frontier equality.  If the closures differed,
observable closure detectability would supply a project with a different
availability-tagged cost, duration, or joint terminal observation, a
contradiction.  Thus the closures and compressed states agree.
\end{proof}

The five componentwise equalities in \eqref{eq:di-definition} make
\(\sim_{\DI}\) an equivalence relation.
\label{prop:di-equivalence-relation}
The proof of \cref{thm:raw-to-compressed-projection} applies verbatim to any
pair \(L\sim_{\DI}L'\): finite-horizon and stationary values agree, the finite
raw-library carrier has a finite quotient, and value is well defined on each
class.  Equality of compressed states implies this equivalence by the forward
characterization just proved.
\label{thm:di-quotient-sufficiency}
Finally, if a representation \(r\) has the property that \(r(L)=r(L')\)
forces equality of all five availability-tagged observations, then equality
of \(r\) implies \(L\sim_{\DI}L'\).
\label{prop:di-refinement-minimality}
This restricted refinement is definitional; it is not a generic
coarsest-bisimulation or full-abstraction result.
